%% =====================================================================
%%  SUPPLEMENTARY INFORMATION
%%  Confidence-Guided Adaptive Reconstruction for Windowed QPE
%%  Companion to B_study.tex. Standalone document.
%%  Build:  pdflatex supplementary / bibtex supplementary / pdflatex x2
%% =====================================================================
\documentclass[preprint,11pt]{elsarticle}

\usepackage{amsmath,amssymb,amsthm}
\usepackage{booktabs}
\usepackage{graphicx}
\usepackage{xcolor}
\usepackage{array}
\usepackage[hidelinks]{hyperref}

\newcommand{\PLACE}[1]{\textcolor{blue}{\textbf{[PLACEHOLDER: #1]}}}

% evidentiary labels, identical to the main article
\newcommand{\PROVEN}{\textsc{(proven)}}
\newcommand{\HEUR}{\textsc{(heuristic)}}
\newcommand{\EMPIR}{\textsc{(empirical)}}
\newcommand{\INVAR}{\textsc{(implementation invariant)}}

\newcommand{\Cw}{C_w}
\newcommand{\dcov}{\delta_{\mathrm{cov}}}
\newcommand{\Smax}{S_{\max}}
\newcommand{\beammax}{\texttt{beam\_max}}

% S-prefixed numbering so nothing collides with the main article
\renewcommand{\thesection}{S\arabic{section}}
\renewcommand{\thetable}{S\arabic{table}}
\renewcommand{\thefigure}{S\arabic{figure}}
\renewcommand{\theequation}{S\arabic{equation}}

\journal{MethodsX}

\begin{document}

\begin{frontmatter}
\title{Supplementary Information for\\
\emph{Where Windowed Shor's Algorithm Actually Fails, and What Fixing
It Buys: A Controlled Study}}
\author[inst1]{Sanjay Lakshmanan R}
\ead{rajasanjay21@gmail.com}
\address[inst1]{Department of Computer Science and Engineering,
Amrita Vishwa Vidyapeetham, Chennai, India}

\begin{abstract}
This document contains the full carry-audit record, the complete
failure-attribution breakdown, the exhaustive statistical and resource
tables, and the artifact map supporting the main article. Nothing here is
required to judge the central claims; everything here is required to
verify them independently. No number in this document differs from the
corresponding number in the main article, and every value is a direct
summary of a released result file (Sec.~\ref{si:artifacts}). The
mathematical derivations of the method's properties, and the bin-level
calibration data, accompany the companion article~\cite{methodpaper}
instead.

Section numbers of the form ``Sec.~S$n$'' refer to this document.
References to the main article are given by section name.
\end{abstract}
\end{frontmatter}

\tableofcontents
\vspace{1em}
% =====================================================================
\section{Failure-study and carry-audit documentation}
\label{si:audit}

\subsection{The corrected carry-check audit, before and after}

The main article records that a carry-related hypothesis (originally
``Module A'') was investigated, that the audit supporting it was found to
have a geometry defect, that the corrected audit returned 0\%
false-acceptance, and that the module was consequently removed rather
than retained without evidence. This section gives the quantitative
record.

Both rows of Table~\ref{tab:audit} are the same 168-combination
exhaustive enumeration over every
$(\mathrm{tail},\mathrm{head},\text{carry-bit})$ triple at overlap widths
1--3. They differ only in whether the synthetic window pair's geometry
was constructed correctly.

\begin{table}[h]
\centering\small
\caption{Carry-check audit, before and after the geometry fix. The
superseded row is retained for auditability. Percentages are of the 168
enumerated combinations.}
\label{tab:audit}
\begin{tabular}{@{}lcc@{}}
\toprule
Outcome & Original (superseded) & Corrected \\
\midrule
True accept  & 19.6\% (33) & 25.0\% (42) \\
True reject  & 19.6\% (33) & 75.0\% (126) \\
False accept & \textbf{55.4\% (93)} & \textbf{0.0\% (0)} \\
False reject & 5.4\% (9)   & 0.0\% (0) \\
\bottomrule
\end{tabular}
\end{table}

\subsection{The defect, precisely}

The audit script constructed its synthetic window pairs with
\texttt{right.start = overlap} instead of
\texttt{right.start = width - overlap}. Because the computed overlap
width is a function of the two windows' start positions and widths, this
pinned the \emph{actual} overlap to exactly one bit regardless of the
nominal \texttt{overlap} value the loop was sweeping. The
$\mathrm{tail}$ and $\mathrm{head}$ bit strings were therefore never
compared at their declared width, and the multi-bit-overlap cases that
the 55.4\% figure was supposed to characterise were never exercised at
all. The swept parameter was disconnected from the geometry it named.

This is worth stating precisely because the failure pattern is reusable:
a synthetic-data generator that accepts a parameter but does not vary the
quantity that parameter names will produce confident, reproducible,
entirely fictitious measurements. The generator was internally
consistent, the loop ran to completion, and the resulting number was
stable across repeats.

\subsection{Three independent checks on the correction}

\begin{enumerate}\itemsep3pt
\item \textbf{Algebraic.} The implementation's acceptance condition and
the audit's ground-truth oracle reduce to the same Boolean function of
the same $(\mathrm{tail},\mathrm{head},c)$ triple
(Lemma 1 of the companion article~\cite{methodpaper}). The original test could therefore only
ever have detected its own arithmetic bug, never a defect in the
predicate it was auditing. This is the strongest of the three checks: it
explains why the corrected result must hold for all inputs, not merely
for the 168 enumerated.
\item \textbf{Empirical re-run.} With the geometry corrected and the same
168-combination enumeration re-executed, the classifier agrees with the
oracle on every combination: 0\% false-accept, 0\% false-reject
(Table~\ref{tab:audit}).
\item \textbf{Instrumented replay.} Across 80 trials reproducing the
overlap-corruption fault-injection condition, a deliberately corrupted
candidate was rejected every time and never survived into a stitched
path. The residual failures under that condition therefore trace to the
beam collapsing to zero surviving paths --- a budget problem, and the
same mechanism modules B, C and D target --- rather than to the check
being fooled.
\end{enumerate}

No change was made to the carry-check implementation; the correction is
entirely to the audit script that measured it. The fix, the regenerated
truth table, and the instrumented replay are all in the released
repository: the corrected audit function in
\texttt{Experiments/phase3\_adversarial.py}, the regenerated truth table
\texttt{phase3\_carry\_truth\_table.csv}, and the instrumented-replay
record. The exact revision will be recorded in the archived snapshot
deposited at acceptance.

\subsection{Coverage of the corrected audit}

The corrected audit's input coverage was additionally checked against
every window geometry the windowing routine can actually produce,
exhaustively over total precisions 2--60 and window sizes 1--16. This
confirms that the degenerate case in which an overlap consumes an entire
window --- the one geometry for which the compatibility check's
precondition $o>0$ could fail to be meaningful --- never arises in
practice.

\subsection{Failure attribution: the full category breakdown}

The attribution figure in the main article groups the 94 failures
observed anywhere in the characterisation campaign by the first rung of a
fixed counterfactual ladder (exact bits $\to$ widened beam $\to$
exhaustive candidates $\to$ noise disabled $\to$ ten-fold shots) that
recovers them. Table~\ref{tab:attrib} gives the counts.

\begin{table}[h]
\centering\small
\caption{Root-cause attribution of all 94 naturally observed failures,
from \texttt{phase5\_attribution.csv}. The first two rows are the
classical-budget mechanisms that motivated the method.}
\label{tab:attrib}
\resizebox{\linewidth}{!}{%
\begin{tabular}{@{}lccl@{}}
\toprule
Category & Count & Share & Mechanism \\
\midrule
Shot limitation            & 48 & 51\% & resolved by additional shots alone \\
Compounded budget interaction & 20 & 21\% & needs candidates \emph{and} beam together \\
\midrule
Noise sensitivity          & 25 & 27\% & resolved by disabling noise \\
Attribution timed out      &  1 &  1\% & ladder did not complete \\
\midrule
\textbf{Total}             & \textbf{94} & & \textbf{72\% classical budget starvation} \\
\bottomrule
\end{tabular}
}
\end{table}

The 20 compounded-interaction failures are the ones invisible to any
single-axis relaxation, and 15 of the 20 occurred at overlap $o=3$. They
are the reason module C is retained in the framework despite its weak
independent evidence: they establish that a coupling between candidate
budget and beam width exists, even though the ablation cannot show that C
is what exploits it.

% =====================================================================

% =====================================================================
\section{Extended statistical and resource tables}
\label{si:tables}

\subsection{Full per-instance ablation}

Table~\ref{tab:heldout} is the complete per-arm, per-instance record
underlying the held-out validation figure and summary in the main
article. All four grids completed 144 seed-paired trials.

\begin{table}[h]
\centering\small
\caption{Complete per-instance ablation. The canonical instance is shown
for reference; the three below it were selected by a rule fixed before
any trial of that stage ran. All grids are 144 seed-paired trials.
$b$ = trials the arm broke (baseline succeeded, arm failed);
$c$ = trials the arm fixed. $p$ is the exact two-sided binomial McNemar
value. Values for the three held-out instances are uncorrected: they are
replications, not members of the confirmatory Holm family
(Sec.~\ref{si:tables}). Bootstrap CIs are 95\% percentile intervals on
the paired difference, $B=2000$ resamples at seed 0.}
\label{tab:heldout}
\resizebox{\linewidth}{!}{%
\begin{tabular}{@{}llcccc@{}}
\toprule
Instance & Arm & Success & Diff.\ vs.\ Baseline [95\% CI] & $(b,c)$ & Exact $p$ \\
\midrule
\multicolumn{6}{@{}l}{\emph{$N{=}21$, $a{=}19$, $r{=}6$ --- canonical (design grid)}}\\
 & Baseline & 52.1\% & ---                              & ---      & ---  \\
 & B        & 77.1\% & $+25.0$ pp $[+18.1,+33.3]$       & $(2,38)$ & $1.5\times10^{-9}$ \\
 & C        & 56.9\% & $\phantom{0}+4.9$ pp $[+1.4,+8.3]$ & $(0,7)$  & $1.6\times10^{-2}$ \\
 & D        & 76.4\% & $+24.3$ pp $[+16.0,+32.6]$       & $(6,41)$ & $1.8\times10^{-7}$ \\
 & Full     & \textbf{100.0\%} & $\mathbf{+47.9}$ \textbf{pp} $[+39.6,+56.3]$ & $(0,69)$ & $3.4\times10^{-21}$ \\
\midrule
\multicolumn{6}{@{}l}{\emph{$N{=}21$, $a{=}2$, $r{=}6$ --- held-out: canonical order, new base}}\\
 & Baseline & 50.0\% & ---                              & ---      & ---  \\
 & B        & 76.4\% & $+26.4$ pp $[+18.8,+34.0]$       & $(1,39)$ & $7.5\times10^{-11}$ \\
 & C        & 54.9\% & $\phantom{0}+4.9$ pp $[+1.4,+9.0]$ & $(0,7)$  & $1.6\times10^{-2}$ \\
 & D        & 79.2\% & $+29.2$ pp $[+20.1,+38.9]$       & $(8,50)$ & $1.6\times10^{-8}$ \\
 & Full     & \textbf{97.9\%}  & $\mathbf{+47.9}$ \textbf{pp} $[+39.6,+56.3]$ & $(0,69)$ & $3.4\times10^{-21}$ \\
\midrule
\multicolumn{6}{@{}l}{\emph{$N{=}15$, $a{=}2$, $r{=}4$ --- held-out: new modulus, new order}}\\
 & Baseline & 88.9\% & ---                              & ---      & ---  \\
 & B        & 100.0\% & $+11.1$ pp $[+6.3,+16.0]$       & $(0,16)$ & $3.1\times10^{-5}$ \\
 & C        & 88.9\% & $\phantom{0}+0.0$ pp $[+0.0,+0.0]$ & $(0,0)$  & $1.00$ \\
 & D        & 88.9\% & $\phantom{0}+0.0$ pp $[-5.6,+5.6]$ & $(8,8)$  & $1.00$ \\
 & Full     & \textbf{100.0\%} & $+11.1$ pp $[+6.3,+16.0]$ & $(0,16)$ & $3.1\times10^{-5}$ \\
\midrule
\multicolumn{6}{@{}l}{\emph{$N{=}21$, $a{=}8$, $r{=}2$ --- held-out: baseline already at ceiling}}\\
 & Baseline & 100.0\% & ---                             & ---      & ---  \\
 & B        & 100.0\% & $\phantom{0}+0.0$ pp $[+0.0,+0.0]$ & $(0,0)$  & $1.00$ \\
 & C        & 100.0\% & $\phantom{0}+0.0$ pp $[+0.0,+0.0]$ & $(0,0)$  & $1.00$ \\
 & D        & 100.0\% & $\phantom{0}+0.0$ pp $[+0.0,+0.0]$ & $(0,0)$  & $1.00$ \\
 & Full     & 100.0\% & $\phantom{0}+0.0$ pp $[+0.0,+0.0]$ & $(0,0)$  & $1.00$ \\
\bottomrule
\end{tabular}
}
\end{table}


\subsection{Realised resource usage per arm}

Table~\ref{tab:secondary} reports what each arm actually spent on the
canonical grid, aggregated over all 144 trials and, for the per-window
quantities, over all windows within them (the grid spans $W=4$ to $W=29$
windows per trial depending on precision and overlap). These are the
quantities from which the matched-resource controls were constructed.

Two features are worth noting. Module B's mean retained candidate count
is $4.25$ against a baseline $k$ of $1.50$, and its maximum reaches 10:
the coverage rule expands substantially beyond the fixed budget it
replaces on ambiguous windows, which is the entire mechanism of its
effect. And the Full arm's mean $m_w$ of $6.52$ exceeds B's $4.25$ even
though both use the same coverage rule --- module D's additional shots
reveal additional distinct patterns, which the coverage rule then has to
cover. The three modules interact through the counts, not only through
the pipeline.

\begin{table}[h]
\centering\small
\caption{Realised resource usage by arm on the canonical 144-trial grid,
from \texttt{phase6\_ablation\_diagnostics.csv}. $m_w$ and $S_w$ are
averaged over all windows of all trials; $P'$ is a per-trial scalar.}
\label{tab:secondary}
\begin{tabular}{@{}lcccccc@{}}
\toprule
Arm & mean $m_w$ & max $m_w$ & mean $P'$ & max $P'$ & mean $S_w$ & max $S_w$ \\
\midrule
Baseline & 1.50 & 2  & 1.50  & 2  & 9.33  & 16 \\
B        & 4.25 & 10 & 1.50  & 2  & 9.33  & 16 \\
C        & 1.50 & 2  & 21.07 & 60 & 9.33  & 16 \\
D        & 1.50 & 2  & 1.50  & 2  & 54.92 & 64 \\
Full     & 6.52 & 13 & 21.44 & 60 & 54.92 & 64 \\
\bottomrule
\end{tabular}
\end{table}

\subsection{Exactness of the matched-resource controls}

The beam-matched control is exact: $P'$ is already a trial-level scalar,
and the control receives that trial's value unchanged, so its realised
beam width equals module C's on every one of the 144 trials. The other
two controls approximate a non-uniform per-window allocation by its mean,
and the direction of the residual error matters for interpreting the null
results, so Table~\ref{tab:matchacc} records it.

\begin{table}[h]
\centering\small
\caption{How closely each matched-resource control reproduced the
adaptive arm's realised budget, canonical grid, from
\texttt{phase7\_matched\_resource.csv}. A control that is
\emph{over}-matched makes the corresponding null result harder to obtain,
not easier.}
\label{tab:matchacc}
\resizebox{\linewidth}{!}{%
\begin{tabular}{@{}llccl@{}}
\toprule
Control & Quantity matched & Adaptive arm & Control & Residual \\
\midrule
candidates-matched vs.\ B & mean candidates/window & $4.25$ & $4.29$ & over-matched by $0.9\%$ \\
shots-matched vs.\ D & mean extra shots/window & $43.82$ & $43.70$ & under-matched by $0.26\%$ \\
beam-matched vs.\ C & realised beam width $P'$ & $21.07$ & $21.07$ & exact on every trial \\
\bottomrule
\end{tabular}
}
\end{table}

Neither residual is large enough to change a conclusion. The
candidates-matched control was, if anything, given slightly more than
module B used. The shots-matched control was given $0.26\%$ less than
module D drew --- a shortfall far too small to account for D's $+2.8$
percentage-point point estimate, and in the direction that favours D
rather than the control.

\subsection{Statistical procedure in full}

\paragraph{Paired significance test.} Outcomes are paired binary, so
McNemar's test \cite{mcnemar1947} is the appropriate choice; Dietterich
\cite{dietterich1998} identifies it as the only test with acceptable
type-I error for algorithms evaluated once per unit, and warns explicitly
against substituting a difference-of-two-proportions test. We use the
exact two-sided binomial form throughout,
$p=2\Pr[\mathrm{Bin}(b+c,\tfrac12)\le\min(b,c)]$, rather than the
$\chi^2$ approximation with or without continuity correction. The exact
form is required here because several comparisons have few discordant
pairs --- the C-vs-baseline contrast has $b+c=7$, far below the
$b+c\gtrsim25$ threshold at which the $\chi^2$ approximation is
trustworthy. The discordant counts $(b,c)$ are reported alongside every
$p$-value, since the test is a statement about those counts alone.

\paragraph{Effect-size intervals.} We report 95\% percentile bootstrap
confidence intervals on the paired success-rate difference, with
$B=2000$ resamples at fixed seed 0. The resampling unit is the trial, and
we bootstrap the \emph{per-pair differences} rather than the two arms
separately, which is what preserves the pairing. Point estimates are
computed once on the full sample; interval bounds are the 2.5th and
97.5th percentiles of the resampled distribution \cite{efron1993}.

\paragraph{Multiple comparisons.} Holm's sequentially rejective procedure
\cite{holm1979} is applied across a pre-declared confirmatory family of
nine: the four arm-vs-baseline tests on the canonical grid, the planned
B-vs-C contrast, and the four adaptive-arm-vs-matched-control tests. It
controls the family-wise error rate without an independence assumption
and is uniformly more powerful than Bonferroni. The held-out instances
are analysed identically but are \emph{not} pooled into the family: they
are replications of an already pre-declared comparison rather than new
hypotheses, and their uncorrected $p$-values are reported with
per-instance discordant counts so a reader may apply any correction they
consider appropriate.

\paragraph{The correction applied.} Table~\ref{tab:holm} gives the
sorted $p$-values with their sequential Holm thresholds and decisions.
Rejection proceeds sequentially and halts at the first retention, so
ranks 6--9 are all retained. \textbf{Module C's improvement over baseline
is the boundary case} at rank 6: $p=1.6\times10^{-2}$ against a threshold
of $1.25\times10^{-2}$. Under the smaller five-comparison family covering
the canonical ablation alone, C would have been retained as significant;
we report the family of nine because it is the one the completed
validation actually tests.

\begin{table}[h]
\centering\small
\caption{Holm--Bonferroni correction across the pre-declared confirmatory
family of nine, $\alpha=0.05$. Rejection proceeds sequentially and halts
at the first retention, so ranks 6--9 are all retained. Module C's
improvement over baseline (rank 6) is the boundary case.}
\label{tab:holm}
\resizebox{\linewidth}{!}{%
\begin{tabular}{@{}clccc@{}}
\toprule
Rank $i$ & Comparison & Exact $p$ & Holm threshold $\alpha/(9-i+1)$ & Decision \\
\midrule
1 & Full vs.\ Baseline                 & $3.4\times10^{-21}$ & $0.0056$ & Reject $H_0$ \\
2 & Full vs.\ shots-matched control    & $7.3\times10^{-12}$ & $0.0063$ & Reject $H_0$ \\
3 & B vs.\ Baseline                    & $1.5\times10^{-9}$  & $0.0071$ & Reject $H_0$ \\
4 & D vs.\ Baseline                    & $1.8\times10^{-7}$  & $0.0083$ & Reject $H_0$ \\
5 & B vs.\ C                           & $1.5\times10^{-5}$  & $0.0100$ & Reject $H_0$ \\
\midrule
6 & C vs.\ Baseline                    & $1.6\times10^{-2}$  & $0.0125$ & \textbf{Retain} $H_0$ \\
7 & D vs.\ shots-matched control       & $0.67$              & $0.0167$ & Retain $H_0$ \\
8 & B vs.\ candidates-matched control  & $1.00$              & $0.0250$ & Retain $H_0$ \\
9 & C vs.\ beam-matched control        & $1.00$              & $0.0500$ & Retain $H_0$ \\
\bottomrule
\end{tabular}
}
\end{table}

\paragraph{The caveat that carries through.} Following Dietterich,
McNemar's test speaks to the difference in error proportion on \emph{one}
evaluation set. It does not account for variability across independent
re-runs of the whole campaign, and two repeats per grid cell are not
sufficient to characterise that. No result in the main article or here
should be read as doing so.

% =====================================================================

% =====================================================================
\section{Artifact map and reproducibility detail}
\label{si:artifacts}

\subsection{Script-to-result-to-section mapping}

Every numeric claim in the main article and in this document is a
mechanical summary of one of the files in Table~\ref{tab:artifacts}. All
files are in the released repository under
\texttt{Data/failure\_study/}, and all generating scripts under
\texttt{Experiments/}. The table lists the files this article consumes;
the companion article~\cite{methodpaper} lists the per-window files
supporting its window-level validation, from the same release.

\begin{table}[h]
\centering\small
\caption{Artifact map. ``Used in'' names the section of the main article
(M) or of this document (S) that consumes the file.}
\label{tab:artifacts}
\scriptsize
\begin{tabular}{@{}>{\raggedright\arraybackslash}p{0.31\linewidth}>{\raggedright\arraybackslash}p{0.30\linewidth}>{\raggedright\arraybackslash}p{0.31\linewidth}@{}}
\toprule
Result file & Produced by & Used in \\
\midrule
\texttt{phase1\_\allowbreak{}ofat.csv} & \texttt{phase1\_\allowbreak{}ofat.py} & M: single-parameter sweep \\
\texttt{phase1b\_\allowbreak{}combined\_\allowbreak{}stress.csv} & \texttt{phase1b\_\allowbreak{}combined\_\allowbreak{}stress.py} & M: combined starvation \\
\texttt{phase2\_\allowbreak{}noise\_\allowbreak{}combined.csv}, \texttt{\_\allowbreak{}single.csv} & \texttt{phase2\_\allowbreak{}noise.py} & M: noise interaction \\
\texttt{phase3\_\allowbreak{}adversarial.csv} & \texttt{phase3\_\allowbreak{}adversarial.py} & M: fault injection \\
\texttt{phase3\_\allowbreak{}ambiguity\_\allowbreak{}boundary.csv} & \texttt{phase3\_\allowbreak{}adversarial.py} & M: ambiguity boundary \\
\texttt{phase3\_\allowbreak{}carry\_\allowbreak{}truth\_\allowbreak{}table.csv} & \texttt{phase3\_\allowbreak{}adversarial.py} & S: Table~\ref{tab:audit} \\
\texttt{phase4\_\allowbreak{}transpile\_\allowbreak{}cliff.csv}, \texttt{phase4\_\allowbreak{}scaling\_\allowbreak{}by\_\allowbreak{}*.csv} & \texttt{phase4\_\allowbreak{}scaling.py} & M: compilation cliff \\
\texttt{phase5\_\allowbreak{}attribution.csv} & \texttt{phase5\_\allowbreak{}attribution.py} & M: attribution; S: Table~\ref{tab:attrib} \\
\texttt{phase6\_\allowbreak{}ablation.csv} & \texttt{phase6\_\allowbreak{}ablation.py} & M: ablation \\
\texttt{phase6\_\allowbreak{}ablation\_\allowbreak{}diagnostics.csv} & \texttt{phase7\_\allowbreak{}generalization.py} & M: held-out, baseline check; S: Tables~\ref{tab:secondary}, \ref{tab:matchacc} \\
\texttt{phase7\_\allowbreak{}heldout\_\allowbreak{}*.csv} & \texttt{phase7\_\allowbreak{}generalization.py} & M: held-out; S: Table~\ref{tab:heldout} \\
\texttt{phase7\_\allowbreak{}matched\_\allowbreak{}resource.csv} & \texttt{phase7\_\allowbreak{}generalization.py} & M: matched controls \\
\midrule
\texttt{plots/*.png} & \texttt{paper\_\allowbreak{}figures.py} and the phase scripts (rendering only; run no trial) & M: all figures \\
\bottomrule
\end{tabular}
\end{table}

\subsection{Determinism, resumability, and the analysis pipeline}

\paragraph{Trials are deterministic given their seed.} Each trial's seed
is derived from its grid coordinates and seeds the simulator's sampling,
so re-running any cell reproduces its counts exactly. This is what
licenses the five-arm within-subjects comparison. It was verified rather
than assumed: the instrumented re-run reproduced all 144 canonical
trials, in all five arms, bit-for-bit.

\paragraph{Long runs are resumable.} A completed trial is skipped on
restart by its seed, so an interrupted campaign resumes without altering
any result --- a practical necessity rather than a convenience, since at
the heaviest grid cells a single trial takes tens of minutes and module
D's resampling rounds each involve a real transpile-and-execute measured
at approximately $3.3$\,s.

\paragraph{Per-cell outcome grids.} The complete per-trial records --- 144
rows $\times$ 5 arms for each of the four instances, plus the four
matched-resource control columns, each row carrying its grid coordinates,
seed, and realised diagnostics --- are released as the CSV files in
Table~\ref{tab:artifacts} rather than reproduced here, since they run to
576 rows and every aggregate in either document is a direct summary of
them.

\subsection{A naming trap in the released summary files}

Anyone recomputing the reported tables from the \emph{summary} CSVs under
\texttt{summary\_tables/} should be aware that their two discordant-count
columns, \texttt{baseline\_only\_fixed\_count} and
\texttt{arm\_only\_broke\_count}, hold the reverse of what their names
suggest: the first is $b$ (baseline succeeded, arm failed) and the second
is $c$ (baseline failed, arm succeeded). The per-trial CSVs are
unambiguous and are what every table in the main article and in this
document was built from. We record the discrepancy rather than silently
renaming the columns, so that anyone who has already used them can tell
which convention they were reading.

\bibliographystyle{elsarticle-num}
\bibliography{refs_B}

\end{document}
